\PageHeader
  {Appendix: tablouri și diagonale în C++}
  {Aceleași reguli de rânduri și coloane, dar în sintaxa unui limbaj real.}
  {Pattern-uri 1 din 3}

\begin{examplebox}[Pătrat plin]
\begin{lstlisting}[style=languagequest,language=C++]
for (int linie = 1; linie <= n; linie++) {
    for (int coloana = 1; coloana <= n; coloana++) {
        cout << "#";
    }
    cout << endl;
}
\end{lstlisting}
\end{examplebox}

\begin{examplebox}[Diagonala principală]
\begin{lstlisting}[style=languagequest,language=C++]
for (int linie = 1; linie <= n; linie++) {
    for (int coloana = 1; coloana <= n; coloana++) {
        if (linie == coloana) cout << "*";
        else cout << ".";
    }
    cout << endl;
}
\end{lstlisting}
\end{examplebox}

\clearpage

\PageHeader
  {Appendix: tablouri și diagonale în Python}
  {Python lasă foarte clar la vedere indentarea buclelor imbricate.}
  {Pattern-uri 2 din 3}

\begin{examplebox}[Pătrat plin]
\begin{lstlisting}[style=languagequest,language=Python]
for linie in range(1, n + 1):
    for coloana in range(1, n + 1):
        print("#", end="")
    print()
\end{lstlisting}
\end{examplebox}

\begin{examplebox}[Diagonala principală]
\begin{lstlisting}[style=languagequest,language=Python]
for linie in range(1, n + 1):
    for coloana in range(1, n + 1):
        if linie == coloana:
            print("*", end="")
        else:
            print(".", end="")
    print()
\end{lstlisting}
\end{examplebox}

\clearpage

\PageHeader
  {Appendix: tablouri și diagonale în Java}
  {Java cere aceeași disciplină a acoladelor și a tipurilor, dar logica rămâne neschimbată.}
  {Pattern-uri 3 din 3}

\begin{examplebox}[Pătrat plin]
\begin{lstlisting}[style=languagequest,language=Java]
for (int linie = 1; linie <= n; linie++) {
    for (int coloana = 1; coloana <= n; coloana++) {
        System.out.print("#");
    }
    System.out.println();
}
\end{lstlisting}
\end{examplebox}

\begin{examplebox}[Diagonala principală]
\begin{lstlisting}[style=languagequest,language=Java]
for (int linie = 1; linie <= n; linie++) {
    for (int coloana = 1; coloana <= n; coloana++) {
        if (linie == coloana) System.out.print("*");
        else System.out.print(".");
    }
    System.out.println();
}
\end{lstlisting}
\end{examplebox}

\clearpage
